% Part: first-order-logic
% Chapter: completeness
% Section: identity

\documentclass[../../../include/open-logic-section]{subfiles}

\begin{document}

\olfileid[id]{fol}{com}{ide}
\olsection{Identitas}

\begin{explain}
Konstruksi model suku yang diberikan pada bagian sebelumnya sudah cukup
untuk menetapkan kelengkapan logika orde pertama bagi himpunan~$\Gamma$
yang tidak memuat~$\eq$. Model suku tersebut memenuhi setiap $!A \in
\Gamma^*$ yang tidak memuat~$\eq$ (dan karenanya setiap $!A \in
\Gamma$). Akan tetapi, konstruksi itu tidak berhasil apabila terdapat
$\eq$. Alasannya, $\Gamma^*$ dalam keadaan tersebut dapat memuat
!!a{sentence}~$\eq[t][t']$, sedangkan dalam model suku nilai setiap suku
adalah suku itu sendiri. Jadi, jika $t$ dan $t'$ merupakan suku yang
berbeda, nilai masing-masing dalam model suku---yakni $t$ dan $t'$---juga
berbeda, sehingga $\eq[t][t']$ salah. Namun, kita dapat mengatasi masalah
ini dengan konstruksi yang dikenal sebagai ``pemfaktoran.''
\end{explain}

\begin{defn}
  Misalkan $\Gamma^*$ adalah himpunan kalimat yang konsisten dan
  !!{complete} dalam~$\Lang L$. Kita mendefinisikan relasi $\approx$
  pada himpunan suku tertutup dari~$\Lang L$ dengan
  \[
  t \approx t' \text{\quad jika dan hanya jika \quad} \eq[t][t'] \in \Gamma^*
  \]
\end{defn}

\begin{prop}
\ollabel{prop:approx-equiv}
Relasi $\approx$ mempunyai sifat-sifat berikut:
\begin{enumerate}
\item $\approx$ bersifat refleksif.
\item $\approx$ bersifat simetris.
\item $\approx$ bersifat transitif.
\item Jika $t \approx t'$, $f$ adalah !!a{function}, dan $t_1$, \dots,
  $t_{i-1}$, $t_{i+1}$, \dots, $t_n$ adalah suku-suku tertutup, maka
\[
\Atom{f}{t_1,\dots, t_{i-1}, t, t_{i+1}, \dots, t_n} \approx
\Atom{f}{t_1,\dots, t_{i-1}, t', t_{i+1}, \dots, t_n}.
\]
\item Jika $t \approx t'$, $R$ adalah !!a{predicate}, dan $t_1$, \dots,
  $t_{i-1}$, $t_{i+1}$, \dots, $t_n$ adalah suku-suku tertutup, maka
\begin{multline*}
\Atom{R}{t_1,\dots, t_{i-1}, t, t_{i+1}, \dots, t_n} \in \Gamma^* \text{ jika dan hanya jika } \\
\Atom{R}{t_1,\dots, t_{i-1}, t', t_{i+1}, \dots, t_n} \in \Gamma^*.
\end{multline*}
\end{enumerate}
\end{prop}

\begin{proof}
Karena $\Gamma^*$ konsisten dan !!{complete}, $\eq[t][t'] \in
\Gamma^*$ jika dan hanya jika $\Gamma^* \Proves \eq[t][t']$. Oleh
karena itu, cukup ditunjukkan hal-hal berikut:
\begin{enumerate}
\item $\Gamma^* \Proves \eq[t][t]$ untuk setiap suku tertutup~$t$.
\item Jika $\Gamma^* \Proves \eq[t][t']$, maka $\Gamma^* \Proves \eq[t'][t]$.
\item Jika $\Gamma^* \Proves \eq[t][t']$ dan $\Gamma^* \Proves
  \eq[t'][t'']$, maka $\Gamma^* \Proves \eq[t][t'']$.
\item Jika $\Gamma^* \Proves \eq[t][t']$, maka
\[
\Gamma^* \Proves
\eq[\Atom{f}{t_1,\dots,t_{i-1},t,t_{i+1},\dots,t_n}][\Atom{f}{t_1,\dots,t_{i-1},t',t_{i+1},\dots,t_n}]
\]
untuk setiap !!{function} beraritas-$n$~$f$ dan suku-suku tertutup
$t_1$, \dots, $t_{i-1}$, $t_{i+1}$, \dots,~$t_n$.
\item Jika $\Gamma^* \Proves \eq[t][t']$ dan
$\Gamma^* \Proves
\Atom{R}{t_1,\dots,t_{i-1},t,t_{i+1},\dots,t_n}$, maka
$\Gamma^* \Proves \Atom{R}{t_1,\dots,t_{i-1},t',t_{i+1},\dots,t_n}$
untuk setiap !!{predicate} beraritas-$n$~$R$ dan suku-suku tertutup
$t_1$, \dots, $t_{i-1}$, $t_{i+1}$, \dots,~$t_n$.
\end{enumerate}
\end{proof}

\begin{prob}
Lengkapilah pembuktian~\olref[fol][com][ide]{prop:approx-equiv}.
\end{prob}

\begin{defn}
Andaikan $\Gamma^*$ adalah himpunan konsisten dan !!{complete} dalam
suatu bahasa~$\Lang L$, $t$ adalah suku tertutup, dan $\approx$ adalah
seperti pada definisi sebelumnya. Maka:
\[
\equivrep{t}{\approx} = \Setabs{t'}{t'\in \Trm[L],\ t' \text{ tertutup},\ t \approx t'}
\]
dan
$\equivclass{\Setabs{t}{t \in \Trm[L],\ t \text{ tertutup}}}{\approx}
= \Setabs{\equivrep{t}{\approx}}{t \in \Trm[L],\ t \text{ tertutup}}$.
\end{defn}

\begin{defn}
\ollabel{defn:term-model-factor}
Misalkan $\Struct M = \Struct M(\Gamma^*)$ adalah model suku
bagi~$\Gamma^*$ dari \olref[mod]{defn:termmodel}. Maka
$\Struct{\equivclass{M}{\approx}}$ ialah !!{structure} berikut:
\begin{enumerate}
\item $\Domain{\equivclass{M}{\approx}} =
  \equivclass{\Setabs{t}{t \in \Trm[L],\ t \text{ tertutup}}}{\approx}$.
\item $\Assign{c}{\equivclass{M}{\approx}} = \equivrep{c}{\approx}$.
\item $\Assign{f}{\equivclass{M}{\approx}}(\equivrep{t_1}{\approx}, \dots,
  \equivrep{t_n}{\approx}) = \equivrep{\Atom{f}{t_1,\dots, t_n}}{\approx}$.
\item $\tuple{\equivrep{t_1}{\approx}, \dots, \equivrep{t_n}{\approx}} \in
  \Assign{R}{\equivclass{M}{\approx}}$ jika dan hanya jika
  $\Sat{M}{\Atom{R}{t_1,\dots, t_n}}$, yakni jika dan hanya jika
  $\Atom{R}{t_1,\dots, t_n} \in \Gamma^*$.
\end{enumerate}
\end{defn}

\begin{explain}
Perhatikan bahwa kita telah mendefinisikan
$\Assign{f}{\equivclass{M}{\approx}}$ dan
$\Assign{R}{\equivclass{M}{\approx}}$ bagi anggota
$\equivclass{\Setabs{t}{t \in \Trm[L],\ t \text{ tertutup}}}{\approx}$
dengan menyebutnya sebagai
$\equivrep{t}{\approx}$, yakni melalui \emph{wakil}~$t \in
\equivrep{t}{\approx}$. Kita harus memastikan bahwa definisi-definisi
ini tidak bergantung pada pemilihan wakil tersebut. Dengan kata lain,
jika pilihan lain~$t'$ menentukan kelas ekuivalensi yang sama
($\equivrep{t}{\approx} = \equivrep{t'}{\approx}$), definisi-definisi itu
harus memberikan hasil yang sama. Sebagai contoh, jika $R$ adalah
!!{predicate} beraritas satu, klausa terakhir dalam definisi mengatakan
bahwa $\equivrep{t}{\approx} \in \Assign{R}{\equivclass{M}{\approx}}$
jika dan hanya jika $\Sat{M}{\Atom{R}{t}}$. Jika, untuk suatu suku lain
$t'$ dengan $t \approx t'$, $\Sat/{M}{\Atom{R}{t'}}$, definisi tersebut
akan mengharuskan $\equivrep{t'}{\approx} \notin
\Assign{R}{\equivclass{M}{\approx}}$. Jika $t \approx t'$, maka
$\equivrep{t}{\approx} = \equivrep{t'}{\approx}$, tetapi tidak mungkin
sekaligus berlaku $\equivrep{t}{\approx} \in
\Assign{R}{\equivclass{M}{\approx}}$ dan $\equivrep{t}{\approx} \notin
\Assign{R}{\equivclass{M}{\approx}}$. Namun,
\olref{prop:approx-equiv} menjamin bahwa hal itu tidak dapat terjadi.
\end{explain}

\begin{prop}
$\Struct{\equivclass{M}{\approx}}$ terdefinisi dengan baik; yakni, jika
$t_1$, \dots, $t_n$, $t_1'$, \dots, $t_n'$ adalah suku-suku tertutup,
dan $t_i \approx t_i'$, maka
\begin{enumerate}
\item $\equivrep{\Atom{f}{t_1,\dots, t_n}}{\approx} =
    \equivrep{\Atom{f}{t_1',\dots, t_n'}}{\approx}$, yakni
  \[
  \Atom{f}{t_1,\dots, t_n} \approx \Atom{f}{t_1',\dots, t_n'}
  \]
  dan
\item $\Sat{M}{\Atom{R}{t_1,\dots, t_n}}$ jika dan hanya jika
  $\Sat{M}{\Atom{R}{t_1',\dots, t_n'}}$, yakni
  \[
    \Atom{R}{t_1,\dots, t_n} \in \Gamma^* \text{ jika dan hanya jika }
    \Atom{R}{t_1',\dots, t_n'} \in \Gamma^*.
  \]
\end{enumerate}
\end{prop}

\begin{proof}
Hasil ini mengikuti \olref{prop:approx-equiv} melalui induksi pada~$n$.
\end{proof}

Seperti dalam kasus model suku, sebelum membuktikan lema kebenaran kita
memerlukan lema berikut.

\begin{lem}
\ollabel{lem:val-in-termmodel-factored} Misalkan $\Struct M = \Struct M
 (\Gamma^*)$. Maka, bagi setiap suku tertutup,
$\Value{t}{\equivclass{M}{\approx}} = \equivrep{t} {\approx}$.
\end{lem}
\begin{proof}
Pembuktiannya serupa dengan pembuktian
\olref[mod]{lem:val-in-termmodel}.
\end{proof}

\begin{prob}
Lengkapilah pembuktian~\olref[fol][com][ide]{lem:val-in-termmodel-factored}.
\end{prob}

\begin{lem}
\ollabel{lem:truth}
$\Sat{\equivclass{M}{\approx}}{!A}$ jika dan hanya jika $!A \in
\Gamma^*$ untuk setiap kalimat~$!A$.
\end{lem}

\begin{proof}
Lakukan induksi pada~$!A$, sama seperti dalam pembuktian
\olref[mod]{lem:truth}. Satu-satunya kasus yang memerlukan perhatian
tambahan ialah ketika $!A \ident \eq[t][t']$.
\begin{align*}
\Sat{\equivclass{M}{\approx}}{\eq[t][t']} & \text{ jika dan hanya jika } \equivrep{t}{\approx} = \equivrep{t'}{\approx}
\text{ (berdasarkan definisi $\Struct{\equivclass{M}{\approx}}$)}\\
& \text{ jika dan hanya jika } t \approx t' \text{ (berdasarkan definisi $\equivrep{t}{\approx}$)}\\
& \text{ jika dan hanya jika } \eq[t][t'] \in \Gamma^* \text{ (berdasarkan definisi $\approx$).}
\end{align*}
\end{proof}

\begin{digress}
Perhatikan bahwa meskipun $\Struct{M(\Gamma^*)}$ selalu !!{enumerable}
dan tak hingga, $\Struct{\equivclass{M}{\approx}}$ mungkin berhingga,
sebab mungkin hanya terdapat berhingga banyak kelas
$\equivrep{t}{\approx}$. Hal ini memang patut diduga karena $\Gamma$
mungkin memuat !!{sentence} yang mengharuskan setiap !!{structure} tempat
kalimat-kalimat itu benar bersifat berhingga. Sebagai contoh,
$\lforall[x][\lforall[y][x = y]]$ adalah !!{sentence} konsisten, tetapi
hanya dipenuhi dalam !!{structure} dengan !!a{domain} yang memuat tepat
satu !!{element}.
\end{digress}

\end{document}
